\chapter{State of the Art}
\section{Image Quality Assessment in Fetal MRI}

Image quality assessment is an important component of fetal magnetic resonance imaging, as the reliability of acquired images directly influences their suitability for clinical interpretation and subsequent quantitative analysis. Compared with conventional MRI examinations, fetal imaging presents additional challenges arising from unpredictable fetal motion, maternal respiration, and variations in fetal position. These factors can introduce substantial differences in image quality between acquisitions, making it necessary to identify images that contain artifacts or insufficient anatomical information \autocite{malamateniou2013motion}.

Traditionally, the assessment of fetal MRI quality has relied on visual inspection by experienced observers. Images are evaluated according to the visibility of anatomical structures and the presence of motion-related degradation, signal loss, or geometric distortions. Although expert assessment provides valuable information about image usability, it requires manual intervention and may be affected by differences in interpretation between observers. Automated image quality assessment (IQA) approaches have therefore been investigated to provide more consistent and reproducible evaluations.

An important development in this area is FetMRQC, introduced by Sanchez et al. \autocite{sanchez2024fetmrqc}. FetMRQC is an open-source machine learning framework designed for automated quality assessment of fetal brain MRI. The method extracts multiple image-quality metrics from unprocessed MRI images and combines these features using a Random Forest model to predict expert-assigned quality ratings. Rather than relying on a single quality indicator, the framework considers several measurable image characteristics to estimate overall image quality.

The authors evaluated FetMRQC using more than 1,600 manually rated fetal brain T2-weighted images collected from four clinical centers and 13 MRI scanners. This multi-center evaluation was designed to investigate the robustness of quality predictions under heterogeneous acquisition conditions. The study demonstrated the potential of combining interpretable image-derived features with machine learning to support automated quality assessment across different clinical environments \autocite{sanchez2024fetmrqc}.

However, the applicability of such methods depends on the imaging modality and the quality characteristics being assessed. FetMRQC was developed and evaluated using anatomical T2-weighted fetal brain MRI. Diffusion MRI presents additional challenges because diffusion encoding introduces sensitivity to motion-related signal variations that may not be adequately represented by quality metrics developed for anatomical images. Consequently, the demonstrated performance of FetMRQC cannot be assumed to transfer directly to fetal diffusion MRI without further evaluation.

Automated quality assessment has also been incorporated into acquisition-oriented fetal diffusion MRI workflows. The HERON framework, developed by Aviles Verdera et al. \autocite{avilesverdera2025heron}, combines automatic fetal-brain localization and segmentation with motion-related quality measures to identify diffusion volumes requiring reacquisition. Its quality-control methodology includes information derived from brain position, changes in segmented brain volume, and variations in image intensity. HERON therefore demonstrates an alternative to general image-quality rating prediction by relating motion and signal characteristics to decisions about individual diffusion acquisitions.

These studies illustrate two complementary directions in fetal MRI quality assessment. FetMRQC investigates machine learning-based prediction of expert quality ratings from image-derived features, while HERON integrates automated motion-related quality assessment into a fetal diffusion MRI acquisition workflow. Their different objectives, imaging modalities, and validation settings highlight the importance of selecting quality-assessment methods according to the intended application. This motivates further investigation of artifact classification methods that use learned image representations in fetal diffusion MRI, while accounting for the specific characteristics of diffusion-weighted acquisitions.

\section{Motion and Artifact Detection in Diffusion MRI}

Diffusion magnetic resonance imaging is particularly sensitive to motion because diffusion-weighted signal measurements depend on the application of diffusion-encoding gradients. Movement during image acquisition can introduce additional signal attenuation, spatial inconsistencies, and distortions that affect the reliability of diffusion measurements. In fetal imaging, these challenges are amplified by unpredictable fetal movements, maternal respiration, and the limited ability to control the position of the subject during an examination. Consequently, motion artifacts represent a major obstacle to obtaining reliable fetal diffusion MRI data \autocite{malamateniou2013motion,snoussi2025haitch}.

Motion-related degradation in diffusion MRI can occur both within individual diffusion volumes and between consecutively acquired volumes. Intra-volume motion may produce signal loss, discontinuities between slices, or inconsistencies in the appearance of anatomical structures. Inter-volume motion introduces differences in fetal-brain position and orientation across successive diffusion measurements, complicating their spatial alignment. These effects are particularly problematic when multiple diffusion directions and weighting levels are required for quantitative analysis. Without adequate quality control or correction, corrupted measurements may introduce errors into diffusion-derived parameters and subsequent microstructural analysis \autocite{snoussi2025haitch,avilesverdera2025heron}.

In addition to motion, diffusion MRI is affected by spatial distortions associated with magnetic field inhomogeneities and the characteristics of echo-planar imaging. Susceptibility differences between tissues and surrounding air can contribute to geometric distortions, while fetal movement may alter the relationship between anatomy and the underlying distortion field. These effects complicate fetal diffusion MRI reconstruction because the acquired images may differ both in anatomical position and spatial geometry. Distinguishing motion-related degradation from other acquisition-related variations is therefore an important consideration when designing automated artifact-assessment methods \autocite{snoussi2025haitch}.

Several approaches have been developed to reduce the effects of motion and image distortion during fetal diffusion MRI processing. Retrospective motion-correction methods generally attempt to align acquired images and reconstruct a consistent representation of the fetal brain. Depending on the acquisition and reconstruction strategy, correction may operate at the volume, slice, or voxel level. However, motion correction alone does not guarantee recovery of information affected by severe signal loss. This motivates the incorporation of outlier detection and data-rejection strategies alongside image registration and reconstruction \autocite{malamateniou2013motion,snoussi2025haitch}.

A relevant example is the HAITCH framework proposed by Snoussi et al. \autocite{snoussi2025haitch}. HAITCH addresses distortion and motion correction in fetal multi-shell diffusion MRI through a combination of acquisition design, image registration, outlier detection, and robust reconstruction. The framework incorporates a blip-reversed dual-echo acquisition strategy for distortion correction and employs motion-estimation and adaptive data-weighting procedures during reconstruction. Its reconstruction methodology accounts for inconsistencies across diffusion measurements and reduces the contribution of corrupted observations. Evaluation on fetal diffusion MRI data demonstrated improvements in reconstructed image quality and the reliability of subsequent diffusion analysis.

An important characteristic of HAITCH is that artifact handling forms part of a broader reconstruction framework. Rather than treating artifact detection as an independent binary classification problem, the approach identifies and downweights unreliable measurements during iterative correction. This distinction is relevant because a method that improves reconstruction quality does not necessarily provide the same output as a classifier intended to distinguish acceptable and artifact-corrupted acquisitions.

A complementary direction is represented by HERON, which investigates motion assessment and prospective reacquisition of affected fetal diffusion MRI volumes \autocite{avilesverdera2025heron}. In contrast to exclusively retrospective correction, acquisition-oriented quality assessment can identify potentially corrupted measurements while further data acquisition remains possible. HERON uses automated fetal-brain segmentation together with motion- and signal-related quality information to support these decisions. Its scanner-integrated methodology and experimental validation are discussed in greater detail in Section 2.5.

The existing literature therefore demonstrates two closely related but distinct approaches to addressing fetal diffusion MRI artifacts. Reconstruction-oriented methods seek to reduce the influence of corrupted measurements after acquisition, whereas acquisition-oriented methods aim to identify measurements that may need to be repeated. Both approaches depend on the reliable recognition of motion-related image degradation, but they differ in their processing requirements, outputs, and evaluation objectives.

These developments provide the scientific context for investigating learned artifact classification from fetal diffusion MRI images. In particular, they highlight the importance of understanding how artifact characteristics vary across individual slices and complete diffusion volumes, and how the available image information influences the reliability of automated quality assessment.

\section{Deep Learning-Based Fetal Brain Segmentation}
\section{Deep Learning Approaches for MRI Artifact Classification}
\section{Automated and Real-Time MRI Quality Control}
\section{Limitations of Existing Approaches and Research Gap}
